\subsection{Exact arithmetic in henselized global function fields}
\label{C:henselian-global-arithmetic-subsection}

\begin{theorem}[Proved: ordinary arithmetic in henselized global fields and all finite extensions]
\label{C:henselian-global-field-arithmetic}
Let $K$ be a supplied finite field tower over
$\mathbf F_q(t)$, with $q$ odd, and let $v$ be
a place supplied by one of the finite certificates
of Lemma~\ref{C:global-function-field-arithmetic}.
Put
\[
 A=\mathcal O_{K,v},\qquad
 H=\operatorname{Frac}(A^h).                              \tag{C.3.HA1}
\]
There are ordinary terminating algorithms for
exact finite names, name validity, equality,
field arithmetic and univariate factorization
over $H$ and over every supplied finite tower
$E/H$, including inseparable towers.
Uniformly in such a tower they also compute
the normalized discrete valuation of each
nonzero element, a uniformizer, the finite
residue field with its full table, residue
maps and lifts, and a complete finite system
of representatives with exact arithmetic
for every specified quotient of the valuation
ring by a power of its maximal ideal.
All finite linear-algebra operations are
therefore effective.

Every such $E$ is henselian and defectless,
and its completion $\widehat E/E$ is regular.
In particular, every irreducible polynomial
over $E$ remains irreducible over
$\widehat E$.
These algorithms use only finite exact
input names; arbitrary infinite series
are not input objects.
\end{theorem}
\begin{proof}
\emph{The initial valued field.}
The supplied global place certificate gives
exact valuations, its finite residue field,
a uniformizer $u\in K$, and residue lifts
by Lemma~\ref{C:global-function-field-arithmetic}.
Thus $A$ has decidable finite names: retain
exact $K$-names of valuation at least zero.
It also has finite matrix-kernel algorithms.
Indeed, a nonzero entry of least valuation
divides every other entry, so elementary
row and column operations over $A$ compute
Smith kernels.  Every division performed
is an exact division in $K$ with integral
quotient certified by the valuation.

The exact-access part of
Theorem~\ref{C:effective-coherent-henselization-WG}
now applies to $A^h$ with finite residue.
For clarity, its finite names consist of
pointed standard \'{e}tale neighborhoods:
localizations of finite free monic algebras
over $A$ at polynomials nonzero at a
specified simple residue root.
Finite $A$-matrix kernels compute
annihilator generators.  A named element
vanishes in the pointed henselization
exactly when one of these generators
has nonzero residue there.
Tensor neighborhoods compare names.
This supplies decidable validity and
equality, arithmetic, residues, and
unit inverses over $A^h$.

The ring $A^h$ is a discrete valuation ring;
its residue field is that of $A$, and
its maximal ideal is generated by the
same $u$
\cite[Lemmas 15.46.1 and 15.46.11,
Tag 07QL]{C-StacksCoherentHenselizations}.
For $0\ne a\in A^h$, compute its
valuation as follows.  If its residue
is nonzero, the valuation is zero.
Otherwise enumerate $b\in A^h$ until
$a=ub$, replace $a$ by $b$, and repeat.
Every required division exists because
the maximal ideal is $uA^h$.
The loop stops because a nonzero element
of a discrete valuation ring has finite
valuation.  Fractions $a/b$, with
decidably nonzero $b$, give all finite
names in $H$, with equality by cross
multiplication and valuation by
subtraction.  This proves the initial
effective valued-field access.

\emph{Defectlessness and completions.}
The field in (C.3.HA1), with its induced
valuation, is the henselization of the
valued one-variable function field
$(K,v)$.  The valuation is trivial
on the finite constant field.
Jarden's Lemma~2.2(b), including its
purely inseparable argument, therefore
says that $H$ is defectless
\cite[Lemmas 2.1--2.2]{C-JardenHenselian1994}.
Explicitly, if $F/H$ is finite, the
unique extension of the valuation has
ramification index $e(F/H)$ and residue
degree $f(F/H)$ satisfying
\[
 [F:H]=e(F/H)f(F/H).                                     \tag{C.3.HA2}
\]
Finite extensions of a henselian valued
field are henselian.  Defectlessness
passes to such an extension $E/H$:
for $F/E$ finite, apply (C.3.HA2) to
$F/H$ and $E/H$, and divide using
multiplicativity of ramification
indices, residue degrees, and field
degrees.
Jarden's Lemma~2.1 then gives regularity
of $\widehat E/E$ for every finite
$E/H$.  Its proof applies to arbitrary
finite extensions, not just separable
ones: immediacy of completion and the
fundamental inequality force
\[
 [\widehat E L:\widehat E]=[L:E]
 \quad\text{for every finite }L/E.                     \tag{C.3.HA3}
\]
Consequently $E[X]/(f)$ remains a field
after tensoring with $\widehat E$
whenever $f\in E[X]$ is irreducible.
This proves the last assertion of the
statement without a perfection assumption.

\emph{Uniform access in a finite extension.}
Suppose the claimed access has been
constructed for a finite extension
$F/H$, and let $E/F$ be a supplied
finite field extension of degree $n$.
Its vector-space presentation gives
exact arithmetic, equality, enumeration,
and the field norm by the determinant
of multiplication.
Let $v_F$ have value group $\mathbf Z$.
The unique valuation on $E$ restricting
to $v_F$ is the computable rational-valued
function
\[
 w(x)=\frac{v_F(N_{E/F}(x))}{n}
       \quad (0\ne x\in E),\qquad w(0)=+\infty .        \tag{C.3.HA4}
\]
To justify this formula in every
characteristic, extend the henselian
valuation to an algebraic closure.
All conjugates of $x$, counted with
their inseparable multiplicities,
have the same value.  Taking the
valuation of their product gives
(C.3.HA4).  Equivalently it follows
from the determinant definition of
the norm and the unique extended
valuation.  No separability is used.

Write $\kappa_F$ for the already
computed finite residue field.
Enumerate positive integers $e,f$
with $ef=n$, monic irreducible
$m(T)\in\kappa_F[T]$ of degree $f$,
and elements $\pi,\theta\in E$.
Use the previously chosen residue
lifts to obtain a monic lift
$\widetilde m(T)$ over the valuation
ring of $F$.
Test the finite conditions
\[
 w(\pi)=1/e,\qquad w(\theta)\geq0,
       \qquad w(\widetilde m(\theta))>0.               \tag{C.3.HA5}
\]
All tests are decidable by (C.3.HA4).
At least one candidate succeeds:
take the actual ramification index
and residue degree, a uniformizer,
and any lift of a generator of the
finite residue extension.
Equation (C.3.HA2) gives $ef=n$.

Every successful candidate is correct.
If the actual indices are $e_0,f_0$,
the first equality of (C.3.HA5)
forces $e\mid e_0$.
The residual root of the irreducible
$m$ gives an embedded degree-$f$
extension of $\kappa_F$ in $\kappa_E$,
so $f\mid f_0$.
Since $ef=n=e_0f_0$, both indices
are exact.  Thus
\[
 v_E=e w,\qquad
 \kappa_E=\kappa_F[T]/(m),\qquad
 \overline\theta=T\bmod m                              \tag{C.3.HA6}
\]
give the normalized valuation, a
uniformizer $\pi$, and the entire
residue field, with a full finite
multiplication table.

Lift a residue element by evaluating
its degree-less-than-$f$ polynomial
in $\theta$, using the fixed lifts
from $\kappa_F$.
For any integral $x\in E$, enumerate
these finitely many lifts $c$ and
select the unique one satisfying
$v_E(x-c)>0$.  This computes its
residue.  These lifts need not form
a coefficient subfield; the residue
calculation uses only their specified
images.
For each requested $N$, the finite
set
\[
 \left\{\sum_{j=0}^{N-1}c_j\pi^j:
                 c_j\text{ is a chosen residue lift}\right\}
                                                               \tag{C.3.HA7}
\]
is a complete irredundant system of
representatives modulo $\pi^N$.
Repeated residue extraction, subtraction,
and division by $\pi$ compute reductions
and all its arithmetic.
This completes effective valuation
and residue access in each finite step.

\emph{Factorization and tower validity.}
It remains to give the terminating
factorization procedure used to
validate the defining irreducible
polynomials in successive towers.
It is uniform in any field $E$
for which the preceding access
has been constructed.
The zero polynomial is recognized and
reported separately.  A nonzero constant
has no irreducible factors; a degree-one
polynomial is returned after its leading
coefficient is removed.  For the remaining
cases, record the nonzero leading
coefficient and make the polynomial monic.
For degree $d$, replace $f(X)$ by
\[
             F(Y)=\pi^{sd}f(Y/\pi^s)                    \tag{C.3.HA8}
\]
with an effectively chosen $s\geq0$
large enough that every coefficient
is integral.  This change preserves
all factor degrees and can be undone.

Run two searches in parallel.
The first enumerates monic proper
factor pairs $G,H\in E[Y]$ and
tests $F=GH$ exactly.
The second, for $N=1,2,\ldots$,
and for every $1\leq r<d$,
enumerates all monic degree-$r$
and degree-$(d-r)$ polynomial
pairs over the finite quotient
in (C.3.HA7), testing their
product against $F$ modulo
$\pi^N$.  If no pair occurs
for any proper degree, this
search reports irreducibility.

The latter report is sound.
Every monic factor of a monic
integral polynomial over $E$
or $\widehat E$ has integral
coefficients: its roots are
integral, and the relevant
valuation ring is integrally
closed.  Thus a proper factor
would appear at every finite
quotient level.
If $F$ is irreducible over $E$,
it is irreducible over $\widehat E$
by (C.3.HA3).
For a fixed factor degree, if
there were compatible-factor
candidates at every level,
the finitely branching tree of
all such coefficient tuples
would have an infinite branch.
Its inverse limit would be a
monic proper factorization over
the complete valuation ring,
a contradiction.  Hence some
finite level has no candidate
for that degree.  There are
only finitely many degrees,
so a common finite level
has none.  The negative search
therefore terminates exactly
in the irreducible case.
The positive search terminates
in the reducible case.
Recursing on proper factors
terminates by degree and
gives the full factorization,
including repeated and
inseparable factors.

Begin with $H$, where access
was already constructed.
At each tower step, this
factorization algorithm tests
irreducibility over the
preceding field before the
next extension is accepted.
Then (C.3.HA4)--(C.3.HA7)
supply its valued-field data,
and the argument repeats.
This inductive order is not
circular.  It proves decidable
tower validity and all the
claimed ordinary algorithms.
\end{proof}
